% !TeX root = ../main.tex
\begin{figure}[H]
\centering
\begin{minipage}[t]{.48\linewidth}
\raggedright\textbf{Source (annotations inferred)}
\begin{lstlisting}[xleftmargin=0pt]
(rec run x =>
  if flip[G](0.5) then x
  else run
    (1 - x + gaussian[E](0, 1))
) 0
\end{lstlisting}
\end{minipage}\hfill
\begin{minipage}[t]{.48\linewidth}
\raggedright\textbf{After replacement (simplified)}
\begin{lstlisting}[xleftmargin=0pt]
(rec run x =>
  if flip[G](0.5) then x
  else run (1 - x)
) 0
\end{lstlisting}
\end{minipage}
\par\medskip
Let $v_i$ be the expected output from loop entry $\texttt{x}=i$
in the transformed program.
\[
 \begin{pmatrix}v_0\\v_1\end{pmatrix}
 =\begin{pmatrix}0&\tfrac12\\\tfrac12&0\end{pmatrix}
  \begin{pmatrix}v_0\\v_1\end{pmatrix}
  +\begin{pmatrix}0\\\tfrac12\end{pmatrix}
 \qquad\Longrightarrow\qquad
 \begin{pmatrix}v_0\\v_1\end{pmatrix}
 =\begin{pmatrix}\tfrac13\\\tfrac23\end{pmatrix}.
\]
\caption{A noisy iteration becomes a finite-state loop.
Both programs have expected output $v_0=1/3$.}
\label{fig:intro-exact}
\Description{The source repeatedly replaces x by one minus x plus Gaussian
noise, stopping with probability one half at each iteration. Removing the noise
leaves two possible loop-entry values, zero and one. Their expected outputs
satisfy v0 equals v1 over two, and v1 equals one half plus v0 over two,
giving initial expected output one third.}
\end{figure}
